\chapter{BGFX, Magnum, LLGL, and Diligent}
\label{ch:abstraction-layers-one}
\label{ch:bgfx-backend}

\texttt{BGFX} integrates the third-party \texttt{bkaradzic/bgfx} rendering library (fetched
directly via CMake's \texttt{FetchContent}, Chapter~\ref{ch:backend-architecture}) behind the
same contract every other renderer implements, using bgfx's own native API --- window and
platform initialization, texture creation, sprite draws, frame submission --- rather than
routing through SDL\_Renderer.

The chapter compares four wrappers whose underlying API can be chosen below CNA.  That second
selection axis is part of every bug report: \texttt{BGFX} on Vulkan and \texttt{BGFX} on
OpenGL are one CNA family but not the same native execution path.

\begin{longtable}{p{0.14\textwidth}p{0.25\textwidth}p{0.22\textwidth}p{0.24\textwidth}}
\toprule
Identity & Backend resolution & Shader strategy & RT / MRT / query \\
\midrule
\endfirsthead
\toprule
Identity & Backend resolution & Shader strategy & RT / MRT / query \\
\midrule
\endhead
\texttt{BGFX} & runtime library selection, CNA override available & bgfx shader assets & yes / yes / device bit \\
\texttt{MAGNUM} & compile-time desktop GL path & raw GLSL compiled at runtime & yes / runtime maximum / yes \\
\texttt{LLGL} & runtime module selection, Linux currently OpenGL & generated SPIR-V and verbatim GLSL & yes / yes / yes \\
\texttt{DILIGENT} & runtime ordered device fallback & one HLSL source cross-compiled at runtime & yes / yes / device features \\
\bottomrule
\end{longtable}

\begin{sourcenote}
The abstraction boundary also owns this renderer's default attachments. CNA passes bgfx a
window, dimensions, renderer selection, and reset flags, but neither the requested
\cnaclass{SurfaceFormat} nor \cnaclass{DepthFormat}. The exact platform backbuffer and
depth/stencil format are therefore bgfx choices which CNA neither selects nor queries, even
though it submits depth/stencil clear flags. A shared SDL fullscreen call can still resize or
restyle the window, after which \texttt{EnsureViewState()} notices the dimensions and calls
\texttt{bgfx::reset}; the empty presentation-format hook does not prevent that window route.
\end{sourcenote}

\section{Clear owns an ordered, aspect-selective view}

BGFX exposes clearing through view state rather than an issue-ordered clear command, so CNA
now makes each public clear its own view. \texttt{RecordClear()} allocates the next per-frame
view, binds the current backbuffer/2D/cube/MRT destination, installs exactly the requested
\texttt{BGFX\_CLEAR\_COLOR}, \texttt{DEPTH}, and/or \texttt{STENCIL} mask and values, then
touches the view so a clear-only cycle is observable. A later draw or clear receives another
view; \texttt{bgfx::setViewOrder()} remaps those views into public call order. Clear therefore
ignores a custom viewport and covers the full target, while draw--clear, clear--draw, repeated
clear, target-switch, cube-face, and MRT sequences retain their chronology. The shared
clear-options and ordered-clear suites discriminate masks and order rather than merely proving
that a final colour appeared; see \S\ref{sec:backend-clear-matrix}.

BGFX's explicit-view model implements target usage without a native load-policy object. The
shared layer maps only \texttt{DiscardContents} to a deterministic black/max-depth/zero-stencil
clear on every bind; that request becomes its own ordered clear view. Preserve and Platform
issue no bind clear, so the attachment survives by construction under \texttt{CLEAR\_NONE}.
The same rule now reaches 2D targets and cube faces, including per-face multisample colour;
depth/stencil follows it too. The cube-usage, per-face-MSAA, and depth/stencil-usage pixel suites
distinguish Discard, Preserve, and Platform rather than relying on the older no-crash usage
smoke test. The full comparison is \S\ref{sec:backend-rendertarget-usage-matrix}.

\section{Recorded BGFX campaign and the Phase 72 closure}

The recorded BGFX validation campaign reported 4{,}375 passes in a 4{,}377-case configured
population, with two hardware skips; its renderer-specific \texttt{ctest} population passed
103 of 105.  These are campaign populations, not a universal current test count; see
Chapter~\ref{ch:counting-discipline}. The two remaining
failures are both environment ceilings, not code defects: one render-target MSAA-resolve test
fails because this project's own headless \texttt{Xvfb} display has no DRI3 support to
resolve against, and one render-target-cube depth-format test (a
\texttt{Depth24Stencil8}-attached cube face producing no color output) has been investigated
three separate times without finding a root cause, and remains formally deferred rather than
silently dropped. A first complete, row-by-row triage of this renderer originally found
thirty-eight real gaps; thirty-seven were fixed in a single closing session, leaving only the
depth-bias item covered below --- which was itself subsequently fixed too, closing the phase
in full.

\section{Choosing bgfx's underlying native renderer}
\label{sec:bgfx-renderer-selection}

Unlike every other renderer in this book, \texttt{BGFX} is itself an abstraction over several
further, distinct native graphics APIs --- OpenGL, OpenGL ES, Vulkan, Direct3D~11/12, and
Metal --- and CNA exposes a real, previously-undocumented mechanism to choose among them: the
\texttt{CNA\_BGFX\_RENDERER} environment variable, resolved through three small, dedicated
functions in \texttt{BgfxRendererSelection.cpp}. \texttt{GetDefaultRendererType()} picks
\texttt{bgfx::RendererType::OpenGL} on Linux specifically and \texttt{Count} (bgfx's own
``let the library auto-detect'' sentinel) everywhere else; \texttt{ParseRendererTypeOverride()}
case-insensitively accepts \texttt{AUTO}, \texttt{OPENGL}, \texttt{OPENGLES}, \texttt{VULKAN},
\texttt{METAL}, \texttt{DIRECT3D11} / \texttt{D3D11}, \texttt{DIRECT3D12} / \texttt{D3D12}, and
\texttt{NOOP} (bgfx's own headless no-op renderer, useful for pure logic testing with no GPU
work issued at all), throwing a \texttt{std::runtime\_error} naming every valid value for
anything else; \texttt{ResolveRendererType()} ties the two together, falling back to the
platform default whenever the environment variable is unset or empty.

\subsection{A worked example: routing one specific test around a real bgfx/GL limitation}

This mechanism exists for a concrete, real reason, not as a speculative escape hatch:
\texttt{modules/renderers/bgfx/examples/\allowbreak{}bgfx\_rendertarget2d\_msaa\_test.cpp}'s own header comment documents a genuine
investigation. In this project's own Linux development sandbox, bgfx's default OpenGL renderer
negotiates only a legacy OpenGL~2.1 context --- confirmed via \texttt{glxinfo} that the
underlying Mesa/RadeonSI driver actually supports OpenGL~4.6 with full hardware acceleration,
so this is bgfx itself requesting an old context by default, not a driver ceiling --- and
MSAA-flagged framebuffer-attached textures do not actually resolve with real sub-pixel blending
under that legacy GL~2.1 path, despite \texttt{bgfx::getCaps()} itself reporting
\texttt{BGFX\_CAPS\_FORMAT\_TEXTURE\_FRAMEBUFFER\_MSAA} support for the exact format requested.
The investigation ruled out a CNA-side defect specifically: forcing bgfx onto its Vulkan
renderer instead (\texttt{CNA\_BGFX\_RENDERER=VULKAN}, routing through the identical GPU) made
both of the test's own checks pass cleanly, including a genuine intermediate (blended) pixel
value along the test triangle's diagonal edge, with \emph{zero} code changes on either side ---
proof the C++ wiring itself is correct and renderer-agnostic, and the gap is specifically in
bgfx's own default legacy-GL-context resolve path. The real \texttt{ctest} registration
(\texttt{modules/renderers/bgfx/examples/CMakeLists.txt}) reflects this finding directly rather than leaving it as
a comment nobody acts on:

\begin{lstlisting}
cna_register_renderer_test(NAME Bgfx_RenderTarget2D_MsaaResolve
    COMMAND cna_test_bgfx_rendertarget2d_msaa
    TIMEOUT 30
    ENVIRONMENT "SDL_VIDEODRIVER=x11;DISPLAY=${CNA_TEST_DISPLAY};CNA_BGFX_RENDERER=VULKAN")
\end{lstlisting}

This is worth reading alongside, not instead of, this chapter's own ``Test baseline'' section
above and \texttt{docs/graphics-renderer-feature-matrix.md}'s repeatedly-reconfirmed count
(spanning multiple independently-dated task closures) that \texttt{Bgfx\_RenderTarget2D\_%
MsaaResolve} remains one of exactly two still-failing tests in \emph{that} project's own
tracked baseline, attributed there to a related but distinctly-framed cause: ``this sandbox's
\texttt{Xvfb} has no DRI3 support.'' The two explanations are not actually in tension --- DRI3
is specifically what a windowed OpenGL context needs to negotiate real hardware acceleration
under X11 indirect rendering, so its absence is a plausible reason bgfx's OpenGL path falls
back to a legacy context in the first place, while Vulkan's own presentation integration does
not depend on the same DRI3/GLX negotiation at all. What this book cannot independently confirm
is whether the specific sandbox behind that repeatedly-cited count genuinely has real Vulkan/GPU
access available to the \texttt{CNA\_BGFX\_RENDERER=VULKAN} override the way the test's own
comment describes, or whether that particular environment lacks real GPU access down either
path --- this book's own screenshot-feasibility investigation (Chapter~\ref{ch:graphicsdevice})
independently confirmed no GPU passthrough at all in \emph{this} book-writing session's own
container, a genuinely different environment from CNA's own separate development sandbox. Both
real, sourced explanations are given here rather than silently picking one, consistent with
this book's own methodology of not resolving an ambiguity the underlying sources themselves
leave open.

\section{DepthBias: a real gap in the underlying library, not just CNA}

\cnaclass{RasterizerState.DepthBias} is the one feature in this chapter whose fix required
working around a genuine absence in bgfx itself, confirmed by reading bgfx's own vendored
source directly: its high-level state API has no depth-bias mechanism of any kind --- no
corresponding flag in its own headers, and no \texttt{glPolygonOffset} call anywhere in its
vendored OpenGL renderer. The fix emulates the feature entirely at the CNA level, via
a per-draw vertex-shader Z-offset (a new \texttt{u\_depthBias} uniform threaded through every
3D vertex shader). \texttt{SlopeScaleDepthBias}, by contrast, remains \emph{deliberately}
unimplemented --- a project-owner decision, not an oversight --- because a true
per-fragment, screen-space-slope computation would force every 3D shader off the GPU's
early-Z optimization path, a cost judged not worth paying for this specific, rarely-used
parameter.

\section{Render-target readback crashes: one root cause, five fixes}

Five of six render-target-related \texttt{glReadPixels} / \texttt{Xvfb} crashes were traced to
a single, shared root cause: bgfx processes its internal views in ascending ID order every
frame, which means any render-target view is always the \emph{last} one processed --- and
therefore still GL-bound at the exact moment \texttt{glReadPixels()} fires for a screenshot.
The fix adds a dedicated ``flush'' view, touched immediately before every screenshot
specifically to force the correct view to be current first. Only the render-target-cube
depth-format crash mentioned above has a different, still-unidentified root cause and was
not resolved by this fix.

\section{View IDs: a free-list-backed pool, not a hardcoded slot}
\label{sec:bgfx-view-ids}

The ``flush view'' fix in the previous section (Task~951) is one visible symptom of a broader
architectural fact worth explaining directly: bgfx organizes all drawing into numbered
\emph{views}, an ordering mechanism distinct from render-target binding itself --- every
\texttt{bgfx::submit()} call targets a specific numeric view id, and bgfx processes views in
ascending id order every frame regardless of the order draws were actually submitted in. This
renderer's own render-target support (\texttt{Detail::AllocateRtViewId()}) originally handed out
a single hardcoded view id --- ``the old hardcoded 1'' its own fix comment calls it --- for
every render target, 2D or cube, MRT or single-target alike. Task~910 replaced that with a real,
free-list-backed pool:

\begin{lstlisting}
static bgfx::ViewId AllocateRtViewId()
{
    static constexpr bgfx::ViewId kMaxBgfxViews = kBackbufferFlushViewId;
    auto& pool = RtViewIdFreeList();
    if (!pool.empty()) {
        const bgfx::ViewId id = pool.back();
        pool.pop_back();
        return id;
    }
    static bgfx::ViewId nextId = 1;   // 0 is permanently reserved for the backbuffer
    if (nextId >= kMaxBgfxViews)
        throw std::runtime_error("Bgfx: exhausted view ids (more concurrently-live "
                                  "render targets than bgfx supports views)");
    return nextId++;
}
\end{lstlisting}

Three real numbers are reserved by construction, not by convention alone: view id \texttt{0} is
permanently the backbuffer's own view and is never handed out by this allocator at all;
\texttt{kBackbufferFlushViewId} --- the same reserved, highest-numbered view Task~951's own
flush-view fix touches every frame, precisely \emph{because} it is guaranteed to sort last in
bgfx's own ascending-id processing order --- sits at the top of the id space, one past the
allocator's own exclusive upper bound; and every id strictly between those two bookends is a
real, currently-live render target's own view. A newly-freed id returns to the pool
(\texttt{RtViewIdFreeList()}) for reuse by the \emph{next} render target created, rather than
the counter monotonically climbing forever --- the free-list is specifically what keeps a
long-running game that creates and destroys render targets across many frames (a shadow-map
pool being cycled, a portal-camera render target rebuilt on level load) from eventually
exhausting bgfx's own hard \texttt{BGFX\_CONFIG\_MAX\_VIEWS} ceiling (256 by default) and
throwing \texttt{std::runtime\_error} on a perfectly ordinary render target creation, purely
because ids were never being recycled. A CNA game with a genuinely unbounded number of
\emph{simultaneously live} render targets --- not created-and-destroyed-over-time, but all
open at once --- can still exhaust the pool; the free-list solves churn, not an unbounded
working set, and the thrown exception is the correct, honest response to that specific case
rather than a bug to route around.

\section{Sampling a render target: a wrong-handle-type cast, since fixed}
\label{sec:bgfx-wrong-handle-cast}

Chapter~\ref{ch:vulkan-backend} named this renderer's own, now-resolved
\cnaclass{RenderTargetCube}-sampling bug as ``a genuine wrong-handle-cast, diagnosed precisely''
in contrast to Vulkan's own still-unresolved equivalent --- this section is that precise
diagnosis, including how it was actually closed. Both \texttt{BgfxSpriteBatchRenderer::Draw} and
\cnaclass{BgfxRenderer}'s own \texttt{envMapping} branch (the code path
\cnaclass{EnvironmentMapEffect}, per Chapter~\ref{ch:stock-effects}, reaches) used to
unconditionally \texttt{static\_cast} whatever \cnaclass{ITextureRenderer} / \cnaclass{%
ITextureCubeRenderer} they were handed to the plain, ordinary-texture concrete type.
\cnaclass{RenderTarget2D} and \cnaclass{RenderTargetCube} are each backed by their own,
entirely separate sibling renderer class --- consistent with
Chapter~\ref{ch:backend-architecture}'s own \cnaclass{IRenderTargetRenderer} interface design,
not a shortcut specific to this renderer. Both of those sibling classes'
\emph{first data member} is a framebuffer handle, not a texture handle.

The reason this compiled cleanly and never crashed, rather than failing loudly the moment it was
exercised, is a genuinely instructive piece of type-confusion: bgfx represents both handle kinds
identically at the ABI level, as a bare \texttt{struct \{ uint16\_t idx; \}}. The cast was
structurally valid C\texttt{++} --- both types really do have that exact layout --- so the
compiler had no basis to reject it, and the resulting object did not crash on access, because
\texttt{idx} really was a valid, in-range integer. It was simply the \emph{wrong} integer: an
index into bgfx's framebuffer handle pool, read and used as though it indexed the texture handle
pool instead. The practical consequence was a render target that silently sampled whatever
texture happened to occupy the coincidentally-matching slot in the wrong pool --- confirmed by
direct memory-layout analysis and by new smoke tests added specifically to prove the cast
\emph{did not crash} on either \cnaclass{RenderTarget2D} or \cnaclass{RenderTargetCube}, which is
exactly why ordinary ``does it crash'' testing had not caught this before. Both call sites are
now fixed (Tasks~873/874, closed by Task~907's own follow-up): \texttt{BgfxSpriteBatchRenderer::%
Draw} now reaches a render target's real dimensions through \cnaclass{ITextureRenderer}'s own
virtual \texttt{GetWidth()} / \texttt{GetHeight()} rather than an unsafe cast, and the
\texttt{envMapping} branch resolves a bound cube texture through a real \texttt{dynamic\_cast} to
a shared \texttt{IBgfxCubeSamplable} interface instead --- a render-to-texture technique (a
portal, a security-camera view, a reflection probe rendered as an ordinary 2D texture rather than
through the dedicated cube-face path) that samples its own render target back through
\cnaclass{SpriteBatch} or \cnaclass{EnvironmentMapEffect} on this renderer can now be trusted,
where it previously could not have been even though nothing about running the code signaled that
anything was wrong. This chapter's own closing summary below already reflected this fix
implicitly, by not naming Tasks~873/874 among the renderer's remaining open items --- the
paragraph above exists so that omission reads as a confirmed, explained resolution rather than
as an unexplained gap between two sections of the same chapter.

\section{A real, found-and-fixed rendering bug: silent black meshes}

One bug is worth naming for how invisible it was until specifically tested: this renderer's
graphics-renderer class never overrode the effect-aware indexed-draw entry point at all. Any
indexed draw bound to an \cnaclass{Effect} whose vertex format lacked a \texttt{Color}
attribute --- which describes essentially any \cnaclass{Model} loaded through
\texttt{ContentManager} (Chapter~\ref{ch:models-meshes}) --- silently fell back to the base
class's default implementation, discarding the entire \cnaclass{GpuDrawParams} struct and
reading an unbound color attribute (which GL defaults to black). The practical, visible
consequence: any loaded 3D model rendered as a solid black silhouette on this renderer,
regardless of its real diffuse color, texture, or lighting --- fixed by adding a real
override.

\subsection{A worked example: recognizing the trigger condition before hitting it}

The black-mesh bug's real trigger condition --- a vertex format with no \texttt{Color}
attribute, which is essentially every \cnaclass{Model} loaded from content rather than
hand-authored with vertex colors baked in --- is checkable directly against a real
\cnaclass{VertexDeclaration}, grounded in \texttt{VertexElement}'s real accessor:

\begin{lstlisting}
bool HasColorAttribute(const VertexDeclaration& decl)
{
    for (const VertexElement& element : decl.GetVertexElements()) {
        if (element.getVertexElementUsageProperty() == VertexElementUsage::Color) {
            return true;
        }
    }
    return false;
}
\end{lstlisting}

Before this bug's fix, any effect-aware indexed draw where \texttt{HasColorAttribute} would
have returned \texttt{false} was exactly the condition that silently rendered solid black on
this renderer --- which, per the finding above, covers essentially every ordinary
\cnaclass{Model} draw, since a typical imported mesh carries position/normal/texture-coordinate
data but no per-vertex color. This is precisely why the bug was invisible for so long despite
being this common: nothing about it depended on an unusual asset or an edge-case vertex
layout, only on whether anyone had specifically pixel-tested a textured, unlit 3D model on
this one renderer.

\section[OcclusionQuery: a real implementation, correctness not provable here]{OcclusionQuery: a real implementation whose correctness cannot be fully proven here}

\cnaclass{OcclusionQuery} was previously a pure no-op on this renderer; the fix wires the
query handle into bgfx's own dedicated submit-with-occlusion-query overload, across all
twelve 3D-draw submission call sites --- no separate correlation/tagging machinery was
needed the way Vulkan's fix required, since bgfx submits synchronously. Two honestly-documented
caveats remain: first, pixel/query correctness genuinely \emph{cannot} be established in this
project's own sandbox, because sabotaging the fix back to a pure no-op produces identical
\texttt{IsComplete()} / \texttt{PixelCount()} output to the fixed version --- the sandbox's own
software Mesa driver returns a non-null result even for a query that was never submitted
anywhere at all, a real ceiling of the software renderer itself, not a CNA defect. Second,
bgfx's own official usage example attaches an occlusion measurement to a separate, dedicated
view specifically so the query is not polluted by other geometry sharing the same depth
buffer; CNA's current fix instead shares whichever view the game's own 3D draw already
targets --- a real, acknowledged architecture gap against true scene-depth query correctness,
not yet attempted.

\section{Instancing: exact transport, with a native-profile shader fault}
\label{sec:bgfx-instancing}

\cnaclass{GraphicsDevice}\allowbreak{}\texttt{::DrawInstancedPrimitives()} reaches a real
renderer implementation --- \texttt{BgfxRenderer::DrawInstancedPrimitivesEx()} is a
substantial implementation, not a stub. There is no standalone
\texttt{modules/renderers/bgfx/examples/*instanced*.cpp} binary, but concluding ``untested''
from that filename search would now be wrong. The shared graphics suite contains BGFX-specific
pixel, cache-cardinality, submission-count, exact native-range, transient-allocation, wireframe,
and \texttt{InstanceFrequency} cases. These prove considerably more than a no-throw smoke test
while also exposing a profile-dependent defect described below.

The shared transport is no longer the old pointer-only bridge. A game marks a bound
\cnaclass{VertexBuffer} as per-instance data by giving its
\cnaclass{VertexBufferBinding} an \texttt{InstanceFrequency} greater than zero when calling
\texttt{SetVertexBuffers()}. \texttt{GraphicsDevice} composes every semantically contributing
binding into the same immutable \texttt{GpuDrawParams::vertexStreams} array used by ordinary
draws, preserving slot, renderer pointer, declaration stride, vertex offset, frequency, and
element count by value. The classic BGFX-supported shape is exactly one stream of each input
rate:

\begin{lstlisting}
FillVertexStreamBindings(p, /*foldedOffset=*/0,
                         /*allowLegacyEmptyDeclarationFallback=*/false);
ValidateVertexStreamRanges(p, baseVertex + minVertexIndex, numVertices,
                           "numVertices", std::to_string(numVertices));
ValidateInstanceStreamRanges(p, instanceCount);
ValidateVertexStreamCapability(p);
\end{lstlisting}

BGFX reads the lowest-slot positive-frequency entry as the instance stream and explicitly
rejects more than one per-vertex or more than one per-instance stream. This is why its
\texttt{MultiStreamVertexInput} capability remains false: a declaration cannot be split across
two geometry buffers and two instance buffers cannot contribute independently. The ordinary
one-plus-one instancing shape does not require that capability and remains valid.

Both binding offsets are now real. The geometry stream's \texttt{VertexOffset} is added to
\texttt{baseVertex} through \texttt{bgfx::setVertexBuffer}'s start-vertex term, while the
instance stream begins copying at its own \texttt{VertexOffset}. BGFX exposes no native
instance divisor in \texttt{setInstanceDataBuffer}; CNA therefore expands frequency grouping
when filling the transient instance buffer: destination instance $i$ copies source record
$\textit{VertexOffset}+\lfloor i/\textit{InstanceFrequency}\rfloor$. Frequency one stays a
single bulk copy. Shared validation first proves that every geometry stream covers the
declared vertex window and that each instance stream contains
$1+\lfloor(n-1)/f\rfloor$ source records, so an undersized binding throws before native
submission instead of reading uninitialized transient memory. Chapter~\ref{ch:backend-architecture},
\S\ref{sec:drawex-instancing-matrix}, compares this transport and its proof boundary across
every renderer.

Two silent exits remain, but their current mechanisms are narrower than the former
pointer-only account. If no semantically contributing binding has a positive
\texttt{InstanceFrequency}, \texttt{FirstInstanceStream(params)} returns null and the renderer
draws nothing. A transient-capacity shortage reported by
\texttt{bgfx::getAvailInstanceDataBuffer()} also skips the draw. Invalid vertex or instance
ranges no longer share that outcome: they throw during the validation described above. A game
whose large instanced batch disappears on BGFX should therefore check that one binding really
is per-instance and then consider transient instance-buffer pressure; malformed ranges should
already have produced a diagnostic.

The remaining correctness fault is in the checked-in instancing shader, not this transport.
\texttt{vs\_instanced3d.sc} constructs the world matrix as
\texttt{mat4(i\_data0, i\_data1, i\_data2, i\_data3)}. BGFX's GLSL profile treats those
arguments as columns, but its HLSL/SPIR-V/Metal/WGSL profiles interpret the generated matrix
constructor differently; BGFX supplies \texttt{mtxFromCols()} precisely to normalize this
cross-profile difference. The dedicated test proves per-instance transforms on OpenGL/OpenGL ES
and pins their failure on Vulkan. Thus ``working'' is justified only for the verified GLSL
native profile today. Exact ranges, frequency expansion, and one-submit behavior remain proven
independently of that shader-orientation bug; changing the source to \texttt{mtxFromCols()} and
regenerating \texttt{bgfx\_shaders.hpp} is the distinct remaining repair.

One further real detail is worth stating precisely because it is easy to assume otherwise:
the draw call's own \texttt{world} matrix parameter is explicitly unused on this path (a
commented-out parameter name in the real signature, \texttt{const Matrix\&
/*world*/}), because per-instance world transforms are expected to arrive already baked into
each instance's own per-instance vertex data, not supplied once as a shared uniform the way a
single, non-instanced draw's \texttt{World} property works --- only \texttt{view} and
\texttt{projection} are combined into the one \texttt{vp} uniform every instance in the batch
shares.

\section{Texture3D and TextureCube readback: copy first, then read}

For a long time, \cnaclass{Texture3D} and \cnaclass{TextureCube} had a more serious
limitation on BGFX than an absent convenience overload: their shared renderer interfaces'
default \texttt{GetData()} implementation was a silent no-op. Task~914 closes that gap without
making ordinary sampled textures pay for a readback-only usage flag. BGFX documents
\texttt{BGFX\_TEXTURE\_READ\_BACK} as incompatible with a texture that is also a normal GPU
resource, so marking the game texture itself readback-capable would undermine the very texture
that an effect needs to sample.

Instead, each \texttt{GetData()} call creates a short-lived transfer texture exactly matching
the requested region. For a cube source it is a 2D RGBA8 destination; \texttt{bgfx::blit()}
selects the requested cube face through the source Z coordinate. For a volume source it is a
matching 3D destination, so the requested $x,y,z,w,h,depth$ sub-volume remains explicit on
both sides of the copy. A temporary destination is marked for a blit destination and for
readback; the ordinary source texture keeps its normal sampled form.

This is not synchronous memory access. \texttt{bgfx::readTexture()} returns the frame at which
its bytes become available. The shared \texttt{AdvanceFramesUntil()} helper advances up to four
BGFX frames before the transfer texture is destroyed, so the caller's buffer is not handed back
before the asynchronous copy completes.
If a selected BGFX renderer cannot create the temporary transfer texture, CNA now logs the
missing \texttt{BGFX\_CAPS\_TEXTURE\_BLIT}/\texttt{READ\_BACK} capability instead of reverting
to the old indistinguishable no-op. The project's Xvfb/llvmpipe OpenGL environment was checked
to support both capabilities before this path was adopted.

The same change also finally made mip allocation testable: \texttt{mipMap} is now passed through
to BGFX's \texttt{createTexture3D} and \texttt{createTextureCube} calls instead of being silently
hardcoded false. Four already-renderer-agnostic programs are registered for BGFX rather than being
reimplemented as cosmetic duplicates: a distinct-slice volume round trip, a non-origin cube
partial rectangle, a three-level volume mip round trip, and a per-face cube mip round trip.
Together they distinguish a correct transfer from reading level zero, the wrong cube face, or a
CPU buffer left untouched by the former interface default.

\section{Viewport uses ordered views; scissor remains per draw}

Both public setters first become stored BGFX-side values, but that similarity ends at
submission. An enabled, nonzero scissor is reissued through \texttt{bgfx::setScissor()} before
every 3D submit and every SpriteBatch submit, because bgfx scissor is per draw and otherwise
resets to ``none.'' The Boolean from \cnaclass{RasterizerState} is retained independently from
the rectangle. \texttt{bgfx\_scissor\_test} proves both halves, then disables the test again to
exclude a one-way/stuck-state implementation.

Viewport is view-level state rather than a draw bit, so CNA preserves changes by selecting an
ordered view segment before each 3D draw or SpriteBatch batch. Consecutive work with the same
target, viewport, and --- for sprites --- Begin transform reuses the segment; a changed viewport
or sprite transform allocates the next view, binds the same target with \texttt{CLEAR\_NONE},
and keeps execution order through the frame's explicit view-order table. This works for the
backbuffer and render targets. SpriteBatch keeps the view rectangle full-sized so a public
clear still covers the target, then folds the captured viewport origin into its orthographic
matrix; 3D sets the sub-rectangle directly. The public min/max depth values remain discarded as
an explicit bgfx limitation.

The original single-subregion proof now sits beside \texttt{Bgfx\_Deferred\_Viewport}, which
queues several viewport changes without a flush, and the transform-specific regression that
separates two Begin matrices even when the viewport is unchanged. The analogous deferred
scissor oracle proves that each draw still receives its own captured rectangle and enable bit.
Section~\ref{sec:backend-viewport-scissor-matrix} records that broader current scope.

\section{Where this leaves the renderer}

The original Phase~72 closure is a useful historical snapshot, but not a freeze-frame: the
later readback and mip work above is one example of a real capability that would be invisible if
the old ``only depth bias remains'' shorthand were repeated literally. The still-deliberate
\texttt{SlopeScaleDepthBias} omission, the non-GLSL instancing matrix-constructor defect, and
the two environment-bound render-target checks remain important limitations, while later
feature work should be evaluated from the current renderer plan and its registered tests rather
than from that earlier closure slogan alone.

\section{MAGNUM: a wrapper with no offline shader generator}

\texttt{MAGNUM} first searches for an installed Magnum and otherwise fetches pinned Corrade
and Magnum revisions.  Emscripten is rejected.  Its default desktop path is GLX, with EGL an
explicit CMake option rather than a runtime driver choice.

Its shader boundary is unique in this cluster: GLSL lives as C++ string literals and is
compiled by \texttt{Magnum::GL::Shader} at runtime.  There is no checked-in generated shader
header and no staleness checker to run.  That removes an offline tool dependency while moving
syntax and driver compatibility failures into device startup.

Extended draws and 2D targets are native.  MRT capability follows the discovered maximum
target count rather than an unconditional constant, anisotropy follows the runtime sampler
limit, and occlusion queries are real.  This is the desirable capability shape: report the
device actually created, not the feature ceiling of the wrapper library in the abstract.

\section{LLGL: a broad library behind a deliberately narrower CNA path}

\texttt{LLGL} is fetched at a pinned release unless \texttt{CNA\_LLGL\_ROOT} supplies it.
The renderer has its own module-selection helper, but the available set is platform-bounded:
on Linux CNA presently offers OpenGL, and an explicit request for Vulkan throws instead of
quietly selecting something else.

Two shader forms are checked in.  Vulkan-style sources are compiled to SPIR-V, while
OpenGL-specific sources are preserved as GLSL; both are assembled into a generated header.
The generator has a check mode, so CI can distinguish a source change from a forgotten
regeneration.

The current renderer owns native extended draws, 2D targets, MRT, and occlusion queries.  Its
documentation nevertheless describes a narrower boundary than upstream LLGL.  Library
support is not automatically CNA support: each resource and state path still requires an
adapter, a capability answer, and observable evidence through the selected module.

\section{DILIGENT: one HLSL program, several device types}

\texttt{DILIGENT} fetches DiligentCore and resolves a device at runtime.  The preference order
is D3D12, Vulkan, D3D11, then OpenGL, filtered by what the build contains; on Linux the real
candidate list is Vulkan followed by OpenGL.  Tests that exercise only the first successful
candidate therefore do not prove the fallback path.

Rather than shipping one shader package per device type, the renderer stores a single HLSL
source and asks DiligentCore to cross-compile it at runtime.  The row-major pragma is part of
that contract because the HLSL-to-GLSL converter recognizes the pragma form.  A shader edit
can thus affect several native backends without generating separate checked-in blobs, but
each translation path still needs its own execution evidence.

Native extended draws, 2D targets, and MRT are implemented.  Occlusion-query support follows
device features.  Of the wrappers in this chapter, Diligent is closest to broad CNA parity,
yet its runtime fallback order creates the sharpest test-design warning: one test binary
registered twice is useful only if the device override genuinely forces two different paths.

\section{Wrapper-selection checklist}

For every reproduction, capture the CNA identity, underlying native backend, shader delivery
path, and whether fallback was automatic or forced.  Then verify the specific feature through
pixels or query values.  Construction proves that two abstraction layers agreed on startup;
it does not prove that state and resources survived both translations.
